% !TeX root = tcolorbox.tex
% include file of tcolorbox.tex (manual of the LaTeX package tcolorbox)
\clearpage
\section{程序库 \mylib{poster}}\label{sec:poster}%
\tcbset{external/prefix=external/poster_}%

本程序库的主要用途是支持用 |tcolorbox| 创建单页海报。

\refEnv{tcbposter} 是一个 |tikzpicture|，
其中可使用 \refCom{posterbox} 命令以按列排布的方式放置 |tcolorbox|。
这一基本概念或多或少借鉴自出色的 \refPkg{baposter} 宏包。

\mylib{raster} 程序库（见 \zvref{sec:raster}）也能生成外观相近的结果，
根据实际项目情况，它可能更为合适。
\begin{itemize}
\item \mylib{raster} 程序库采用流式布局概念，就像常规的文本流。
  文本流（文本框流）是一条近乎无尽的带状流，被折成若干行（和段落），
  行再被分到各页。\mylib{raster} 将文本框调整为合适的尺寸，
  以便令人舒适地填满各行与各页。
\item \mylib{poster} 程序库支持在页面内相当自由地放置文本框。
  基本上，文本框的放置方式如同在 |tikzpicture| 中放置 |node|。
  与 \mylib{raster} 不同，这是一个\emph{单}页，而不是页面流。
  海报被划分为若干列和行。系统会在一定程度上引导你使用列（或跨列）
  来定位和调整尺寸，而行的放置则完全是可选的。
\end{itemize}
本程序库的创建源于 Ignasi 的提议。

\begin{marker}
在一个 |tikzpicture| 内部不应再嵌套 |tikzpicture|。
\mylib{poster} 程序库打破了这一规则。请注意：主 |tikzpicture|
与嵌入在各 |tcolorbox| 内部的 |tikzpicture| 之间可能会产生一些
不希望出现的相互影响。
\end{marker}

该程序库可通过宏包选项加载，或在导言区中使用：
\begin{dispListing}
\tcbuselibrary{poster}
\end{dispListing}
这同时会加载以下程序库：
\mylib{skins}（见 \zvref{sec:skins}）、
\mylib{breakable}（见 \zvref{sec:breakable}）、
\mylib{magazine}（见 \zvref{sec:magazine}）以及
\mylib{fitting}（见 \zvref{sec:fitting}）。


%--------------------------
\subsection{概览}\label{subsec:poster_overview}


\begin{tcolorbox}[base example,hyperurl={tcolorbox-tutorial-poster.pdf},title=点击查看教程]
若想全面了解 \mylib{poster} 程序库及其功能，
请参阅 |tcolorbox| 文档中的\textbf{海报教程}，它是 |tcolorbox| 文档的一部分：\par
\texttt{tcolorbox-tutorial-poster.pdf}
\end{tcolorbox}



\clearpage
%--------------------------
\subsection{主海报环境}\label{subsec:poster_environment}

\begin{docEnvironment}[doc new=2017-07-03]{tcbposter}{\oarg{options}}
  该环境创建一个 |tikzpicture|，并根据给定的 \meta{选项} 设置相应的附加参数。
  基本上，使用 \refCom{posterbox} 和 \refEnv{posterboxenv}
  将 |tcolorbox| 作为节点放入该环境，也可以使用额外的 \tikzname\ 代码。
  作为 \meta{选项}，可以使用所有 |/tcb/posterset/| 下的键，即：
\begin{itemize}
\item\refKey{/tcb/posterset/poster}：海报设置，如列、行、尺寸……
\item\refKey{/tcb/posterset/coverage} 与 \refKey{/tcb/posterset/no coverage}：
  用于设置环绕外围的 |tcolorbox| 的背景和边距。
\item\refKey{/tcb/posterset/boxes}：海报中所用各 |tcolorbox| 的样式。
\item\refKey{/tcb/posterset/fontsize}：所用字体的缩放。
\end{itemize}

\begin{exdispExample}{tcbposter}
  \begin{tcbposter}[
    poster = {showframe,height=10cm,spacing=2mm},
    boxes  = {beamer,colframe=blue!50!black,colback=blue!50,colupper=yellow!50},
  ]
  \posterbox{name=A,column=3,row=2}{My first box}
  \posterbox[adjusted title=Second box]
            {name=B,column=2,span=2,below=A}{My second box}
  \posterbox[adjusted title=Third box]
            {name=C,column=2,between=B and bottom}{My third box}
  \end{tcbposter}
\end{exdispExample}
\end{docEnvironment}

\clearpage
  在 \refEnv{tcbposter} 内部，预定义了若干 \tikzname\ 节点。
  这些节点共享同一个 \refKey{/tcb/poster/prefix}，默认为 |TCBPOSTER@|。
  此前缀用于区分海报节点与任何内嵌 |tikzpicture| 环境中的局部节点。
  使用 \refCom{posterbox} 及其放置选项时你永远用不到此前缀，
  但如果你想用纯 \tikzname\ 代码引用某个预定义节点，则会用到它。
  预定义节点（不显示前缀）如下：
  \begin{itemize}
  \item|poster|：定义海报的包围盒（不含 coverage 外围框）。
  \item|top|：顶部位置加行间距
  \item|bottom|：底部位置减行间距
  \item|middle|：垂直方向居中位置
  \item|col1|、|col2|、\ldots：第 1 列、第 2 列……的包围盒
  \item|row1|、|row2|、\ldots：第 1 行、第 2 行……的包围盒
  \end{itemize}
  其余节点通过 \refKey{/tcb/posterloc/name} 选项定义。

  \begin{marker}
  切勿在一个 \refEnv{tcbposter} 内部再使用 \refEnv{tcbposter}。
  但若你执意如此，请为内嵌的海报使用不同的 \refKey{/tcb/poster/prefix}，
  否则必然会一团糟。
  \end{marker}

  \refEnv{tcbposter} 内部还有若干属性，可能对高级代码有用
  （初读时可跳过以下内容）：
  \begin{itemize}
  \item\docAuxCommand{tcbposterwidth}：海报的宽度（不含边距）。
  \item\docAuxCommand{tcbposterheight}：海报的高度（不含边距）。
  \item\docAuxCommand{tcbpostercolspacing}：列间距。
  \item\docAuxCommand{tcbposterrowspacing}：行间距。
  \item\docAuxCommand{tcbpostercolumns}：列数。
  \item\docAuxCommand{tcbposterrows}：行数。
  \item\docAuxCommand{tcbpostercolwidth}：单列宽度。
  \item\docAuxCommand{tcbposterrowheight}：单行高度。
  \end{itemize}

\medskip
\begin{docCommand}[doc new=2017-07-03]{tcbposterset}{\marg{options}}
  为当前 \TeX\ 分组内此后的每个 \refEnv{tcbposter} 设置选项。
  例如，可用以下方式为整篇文档定义行列数：
\begin{dispListing}
\tcbposterset{poster={columns=2,rows=3}}
\end{dispListing}
  所有可用选项见 \refEnv{tcbposter}。
\end{docCommand}


\clearpage
%--------------------------
\subsection{海报设置}\label{subsec:poster_settings}

\begin{postersetTcbKey}[][doc new=2017-07-03]{poster}{=\marg{option list}}{样式，无默认值}
  此选项可在 \refEnv{tcbposter} 内部及 \refCom{tcbposterset} 中使用，
  用于设置给定的海报 \meta{选项列表}，例如：
\begin{dispListing}
\tcbposterset{poster={width=20cm,height=15cm}}
\end{dispListing}
  关于 \meta{选项列表}，参见以下各键。
\end{postersetTcbKey}


\begin{posterTcbKey}[][doc new=2017-07-03]{columns}{=\meta{number}}{无默认值，初始 |3|}
  设置 \refEnv{tcbposter} 的列 \meta{数}。
\begin{exdispExample}{columns}
  \begin{tcbposter}[
    poster = {showframe,columns=5,rows=2,spacing=1mm,height=4cm},
  ]
  \end{tcbposter}
\end{exdispExample}
\end{posterTcbKey}

\begin{posterTcbKey}[][doc new=2017-07-03]{rows}{=\meta{number}}{无默认值，初始 |4|}
  设置 \refEnv{tcbposter} 的行 \meta{数}。
\end{posterTcbKey}


\begin{posterTcbKey}[][doc new=2017-07-03]{colspacing}{=\meta{length}}{无默认值，初始 |4mm|}
  将 \meta{长度} 设为列间距。
\end{posterTcbKey}

\begin{posterTcbKey}[][doc new=2017-07-03]{rowspacing}{=\meta{length}}{无默认值，初始 |4mm|}
  将 \meta{长度} 设为行间距。
\end{posterTcbKey}

\begin{posterTcbKey}[][doc new=2017-07-03]{spacing}{=\meta{length}}{样式，无默认值，初始 |4mm|}
  将 \meta{长度} 设为列与行的间距。
\end{posterTcbKey}


\begin{posterTcbKey}[][doc new=2017-07-03]{showframe}{\colOpt{=true\textbar false}}{默认 |true|，初始 |false|}
  在海报创建过程中显示红色辅助网格，作为视觉参考。
  同时显示每个 \refKey{/tcb/posterloc/name}。
\end{posterTcbKey}


\begin{posterTcbKey}[][doc new=2017-07-03]{width}{=\meta{length}}{无默认值，初始 \cs{linewidth}}
  将 \meta{长度} 设为海报宽度。对于常规海报，无需手动设置此项。
  尤其是，当使用了 \refKey{/tcb/posterset/coverage} 时，
  请改用 |coverage={width=|\meta{长度}|}| 来更改整体宽度。
\end{posterTcbKey}

\enlargethispage*{1cm}

\begin{posterTcbKey}[][doc new=2017-07-03]{height}{=\meta{length}}{无默认值，初始未设置}
  将 \meta{长度} 设为海报高度。对于常规海报，无需手动设置此项，
  而会自动设为合适的值。若使用了 \refKey{/tcb/posterset/coverage}，
  则 |coverage={height=|\meta{长度}|}| 与 |poster={height=|\meta{长度}|}|
  二选一即可。
\end{posterTcbKey}


\begin{posterTcbKey}[][doc new=2017-07-03]{prefix}{=\meta{name}}{无默认值，初始 |TCBPOSTER@|}
  \meta{名称} 被设为前缀，用于 \mylib{poster} 程序库自动生成的
  所有 \tikzname\ 节点。这包括像 |top|、|bottom| 等预定义节点，
  以及用 \refKey{/tcb/posterloc/name} 定义的节点。
  另见 \zvref{subsec:poster_environment}。对于常规海报，此值可保持默认。
\end{posterTcbKey}


%--------------------------
\subsection{外围覆盖框}\label{subsec:poster_coverage}

\begin{postersetTcbKey}[][doc new=2017-07-03]{coverage}{=\marg{option list}}{样式，无默认值}
  此选项可在 \refEnv{tcbposter} 内部及 \refCom{tcbposterset} 中使用，
  为海报添加一个可选的外围框，即一个套用给定 \meta{选项列表} 的
  环绕 |tcolorbox|。在此可设置海报的边距和背景。
  \emph{外围框} 带有若干适合此用途的默认 |tcolorbox| 设置：
\begin{dispListing}
enhanced, frame hidden, sharp corners, boxsep=0pt, boxrule=0pt,
top=4mm, bottom=4mm, left=4mm, right=4mm,
toptitle=2mm, bottomtitle=2mm, colback=white
\end{dispListing}

\meta{选项列表} 可包含任意 |tcolorbox| 选项。

\begin{exdispExample}{coverage}
\begin{tcbposter}[
  poster   = {showframe,spacing=1mm},
  coverage = {height=5cm,
              interior style={top color=yellow,bottom color=yellow!50!red},
              watermark text={My Poster},watermark color=white,
             },
]
\end{tcbposter}
\end{exdispExample}

\begin{itemize}
\item 对于常规海报，使用 \refKey{/tcb/spread} 选项
  可让海报外围框占满整个页面。
\item 海报边距可通过 \refKey{/tcb/left}、\refKey{/tcb/right}、
  \refKey{/tcb/top}、\refKey{/tcb/bottom} 调整。
\item 海报背景可通过 \refKey{/tcb/colback}、
  \refKey{/tcb/interior style}、\refKey{/tcb/interior style image} 等更改。
\item 请勿将 \refKey{/tcb/poster/width} 与 \refKey{/tcb/poster/height}
  与\emph{外围框} 同时使用。注意：你可以在\emph{外围框} 的
  \meta{选项列表} 内使用 \refKey{/tcb/width} 和 \refKey{/tcb/height}。
  另外，当使用了 \refKey{/tcb/spread} 时，这也并非必要。
\end{itemize}
\end{postersetTcbKey}


\begin{postersetTcbKey}[][doc new=2017-07-03]{no coverage}{}{样式，无取值，初始已设置}
  完全移除外围的 |tcolorbox|。
\end{postersetTcbKey}

\clearpage
%--------------------------
\subsection{通用文本框设置}\label{subsec:poster_boxsettings}


\begin{postersetTcbKey}[][doc new=2017-07-03]{boxes}{=\marg{option list}}{样式，无默认值}
  此选项可在 \refEnv{tcbposter} 内部及 \refCom{tcbposterset} 中使用，
  用于设置海报内各 |tcolorbox| 的样式。
  \meta{选项列表} 可包含任意 |tcolorbox| 选项，
  但此处文本框尺寸选项预计用处不大，因为尺寸由放置选项决定。

\begin{exdispExample}{boxes}
\begin{tcbposter}[
  poster   = {spacing=2mm,columns=3,rows=2},
  coverage = {height=5cm,
              interior style={top color=yellow,bottom color=yellow!50!red},
             },
  boxes    = {sharp corners=downhill,arc=3mm,boxrule=1mm,
              colback=white,colframe=cyan,
              title style={left color=black,right color=cyan},
              fonttitle=\bfseries}
]
  \posterbox[adjusted title=First]{column=1,row=1,span=2}{First box}
  \posterbox[adjusted title=Second]{column=1,row=2,span=2}{Second box}
  \posterbox[adjusted title=Third]{column=3,row=1,rowspan=2}{Third box}
\end{tcbposter}
\end{exdispExample}

\end{postersetTcbKey}


%--------------------------
\subsection{字体缩放}\label{subsec:poster_fontsize}

\begin{postersetTcbKey}[][doc new=2017-07-03]{fontsize}{=\meta{length}}{样式，无默认值，初始未设置}
  此选项可在 \refEnv{tcbposter} 内部及 \refCom{tcbposterset} 中使用。
  它借助 \refKey{/tcb/fit basedim} 和 \refKey{/tcb/fit fontsize macros}
  将 |\normalsize| 重定义为 \meta{长度}，并相应地重定义 |\small|、|\large|
  等其他标准字号宏。\par
  这需要像 \refPkg[lm]{lmodern} 这样可自由缩放的字体族才能生效。
  若未使用 \refKey{/tcb/posterset/fontsize}，
  则这些标准字号宏不会有任何改动。

\begin{dispListing}
\begin{tcbposter}[
  poster   = {spacing=2mm,columns=3,rows=2},
  coverage = {height=5cm,
              interior style={top color=yellow,bottom color=yellow!50!red},
             },
  fontsize = 15pt,   % <--- \normalsize is now 15pt
]
...
\end{dispListing}
\end{postersetTcbKey}


\clearpage
%--------------------------
\subsection{文本框放置}\label{subsec:poster_boxplacement}

\begin{docCommand}[doc new=2017-07-03]{posterbox}{\oarg{options}\marg{placement}\marg{box content}}
  在 \refEnv{tcbposter} 环境内部，此命令将一个带有附加
  |tcolorbox| \meta{选项} 和给定 \meta{文本框内容} 的 |tcolorbox|
  放置到由 \meta{放置方式} 决定的位置。
  所有 \meta{放置方式} 选项见下文。
  注意：\meta{文本框内容} 不能包含 \emph{verbatim}（原样抄录）材料，
  见 \refEnv{posterboxenv}。
\begin{exdispExample}{posterbox}
\begin{tcbposter}[
  poster = {showframe,height=4cm,spacing=2mm,rows=2},
  boxes  = {beamer,colframe=blue!50!black,colback=blue!50,colupper=yellow!50},
]
\posterbox[title=My title]{name=A,column=2,row=2}{My first box}
\end{tcbposter}
\end{exdispExample}
\end{docCommand}

\begin{docEnvironment}[doc new=2017-07-03]{posterboxenv}{\oarg{options}\marg{placement}}
  这是 \refCom{posterbox} 的环境版本，即在 \refEnv{tcbposter}
  环境内部，它将一个带有附加 |tcolorbox| \meta{选项} 和给定
  \meta{环境内容} 的 |tcolorbox| 放置到由 \meta{放置方式} 决定的位置。
  与 \refCom{posterbox} 不同，\meta{环境内容} 允许包含 \emph{verbatim} 材料。
  注意：\refCom{posterbox} 的实现比 \refEnv{posterboxenv} 的实现更高效。

\enlargethispage*{1cm}
\begin{exdispExample}{posterboxenv}
\begin{tcbposter}[
  poster = {showframe,height=4cm,spacing=2mm,rows=2},
  boxes  = {size=small,beamer,
            colframe=blue!50!black,colback=blue!50,colupper=yellow!50},
]
\begin{posterboxenv}[title=My title]{name=A,column=2,between=top and bottom}
  My first box.
  \begin{tcblisting}{size=small,colback=yellow!10}
My \textbf{first}
poster listing.
  \end{tcblisting}
\end{posterboxenv}
\end{tcbposter}
\end{exdispExample}

\end{docEnvironment}


\clearpage
\begin{posterlocTcbKey}[][doc new=2017-07-03]{name}{=\meta{name}}{无默认值，初始 |@|}
  将 \meta{名称} 设为当前 \refCom{posterbox} 或 \refEnv{posterboxenv} 的引用名。
  一个 \tikzname\ 形状名会由 \refKey{/tcb/poster/prefix} 与 \meta{名称}
  自动组合而成。
\begin{exdispExample}{name}
\begin{tcbposter}[
  poster = {showframe,height=2.5cm,spacing=2mm,rows=2},
  boxes  = {beamer,colframe=blue!50!black,colback=blue!50,colupper=yellow!50},
]
\posterbox{name=A,column=2,row=2}{My first box}
\node[below right=4mm,fill=yellow] (X) at (TCBPOSTER@poster.north west) {Example A};
\draw[blue,very thick,->] (X) |- (TCBPOSTER@A);
\end{tcbposter}
\end{exdispExample}
\end{posterlocTcbKey}


\begin{posterlocTcbKey}[][doc new=2017-07-03]{column}{=\meta{number}}{无默认值，初始 |1|}
  将文本框放置到 \meta{数字} 所指定的列。
  若 \refKey{/tcb/posterloc/span} 不为 |1|，
  则文本框沿第 \meta{数字} 列的左侧对齐。
\begin{exdispExample}{column}
\begin{tcbposter}[
  poster = {showframe,height=2.5cm,spacing=2mm,rows=2},
  boxes  = {beamer,colframe=blue!50!black,colback=blue!50,colupper=yellow!50},
]
\posterbox{row=1,column=2,span=2}{First box}
\posterbox{row=2,column=2,span=0.8}{Second box}
\end{tcbposter}
\end{exdispExample}
\end{posterlocTcbKey}

\enlargethispage*{1cm}
\begin{posterlocTcbKey}[][doc new=2017-07-03]{column*}{=\meta{number}}{无默认值，初始未设置}
  将文本框放置到 \meta{数字} 所指定的列。
  若 \refKey{/tcb/posterloc/span} 不为 |1|，
  则文本框沿第 \meta{数字} 列的右侧对齐。
\begin{exdispExample}{columnstar}
\begin{tcbposter}[
  poster = {showframe,height=2.5cm,spacing=2mm,rows=2},
  boxes  = {beamer,colframe=blue!50!black,colback=blue!50,colupper=yellow!50},
]
\posterbox{row=1,column*=2,span=2}{First box}
\posterbox{row=2,column*=2,span=0.8}{Second box}
\end{tcbposter}
\end{exdispExample}
\end{posterlocTcbKey}


\clearpage
\begin{posterlocTcbKey}[][doc new=2017-07-03]{span}{=\meta{number}}{无默认值，初始 |1|}
  将当前文本框的宽度设为跨越 \meta{数字} 列。
  \meta{数字} 也可取实数，如 |0.5| 或 |1.7|。
  示例见 \refKey{/tcb/posterloc/column} 与 \refKey{/tcb/posterloc/column*}。
\end{posterlocTcbKey}

\begin{posterlocTcbKey}[][doc new=2017-07-03]{row}{=\meta{number}}{无默认值，初始未设置}
  若使用此选项，文本框将放置到 \meta{数字} 所指定的行。
  同时，高度会按 \refKey{/tcb/posterloc/rowspan} 设为固定值。
\begin{exdispExample}{row}
\begin{tcbposter}[
  poster = {showframe,height=2.5cm,spacing=2mm,rows=2},
  boxes  = {beamer,colframe=blue!50!black,colback=blue!50,colupper=yellow!50},
]
\posterbox{row=1,column=1}{First box}
\posterbox{row=1,column=2,rowspan=2}{Second box}
\posterbox[natural height]{row=1,column=3}{Third box}
\end{tcbposter}
\end{exdispExample}
\end{posterlocTcbKey}

\begin{posterlocTcbKey}[][doc new=2017-07-03]{rowspan}{=\meta{number}}{无默认值，初始 |1|}
  将当前文本框的高度设为跨越 \meta{数字} 行。
  \meta{数字} 也可取实数，如 |0.5| 或 |1.7|。
\begin{exdispExample}{rowspan}
\begin{tcbposter}[
  poster = {showframe,height=2.5cm,spacing=2mm,rows=2},
  boxes  = {beamer,colframe=blue!50!black,colback=blue!50,colupper=yellow!50},
]
\posterbox{row=1,column=1,rowspan=0.9}{First box}
\posterbox{row=1,column=2,rowspan=1.5}{Second box}
\posterbox{row=1,column=3,rowspan=2}{Third box}
\end{tcbposter}
\end{exdispExample}
\end{posterlocTcbKey}

\begin{posterlocTcbKey}[][doc new=2017-07-03]{fixed height}{}{无取值，初始 |0pt|}
  将当前文本框的高度设为跨越 \refKey{/tcb/posterloc/rowspan}
  所指定的行数。当未使用 \refKey{/tcb/posterloc/row}
  而使用了其他高度放置选项时，可使用此选项。
\end{posterlocTcbKey}


\clearpage
\begin{posterlocTcbKey}[][doc new=2017-07-03]{below}{=\meta{name}}{无默认值，初始 |top|}
  将文本框放置到另一个名为给定 \meta{名称} 的文本框下方。
  \meta{名称} 也可以是预定义节点，见 \zvref{subsec:poster_environment}。
\begin{exdispExample}{below}
\begin{tcbposter}[
  poster = {showframe,height=3cm,spacing=2mm,rows=2},
  boxes  = {beamer,colframe=blue!50!black,colback=blue!50,colupper=yellow!50},
]
\posterbox{name=A,column=1,below=top}{First box}
\posterbox{name=B,column=1,below=A}{Second box}
\posterbox{name=C,column=2,below=B}{Third box}
\posterbox{name=D,column=3,below=row1}{Fourth box}
\end{tcbposter}
\end{exdispExample}
\end{posterlocTcbKey}


\begin{posterlocTcbKey}[][doc new=2017-07-03]{above}{=\meta{name}}{无默认值，初始未设置}
  将文本框放置到另一个名为给定 \meta{名称} 的文本框上方。
  \meta{名称} 也可以是预定义节点，见 \zvref{subsec:poster_environment}。
\begin{exdispExample}{above}
\begin{tcbposter}[
  poster = {showframe,height=3cm,spacing=2mm,rows=2},
  boxes  = {beamer,colframe=blue!50!black,colback=blue!50,colupper=yellow!50},
]
\posterbox{name=A,column=1,above=bottom}{First box}
\posterbox{name=B,column=1,above=A}{Second box}
\posterbox{name=C,column=2,above=B}{Third box}
\posterbox{name=D,column=3,above=row2}{Fourth box}
\end{tcbposter}
\end{exdispExample}
\end{posterlocTcbKey}


\clearpage
\begin{posterlocTcbKey}[][doc new=2017-07-03]{at}{=\meta{name}}{无默认值，初始未设置}
  将文本框放置到名为给定 \meta{名称} 的位置。
  这很可能是一个预定义节点，见 \zvref{subsec:poster_environment}。
\begin{exdispExample}{at}
\begin{tcbposter}[
  poster = {showframe,height=3cm,spacing=2mm,rows=2},
  boxes  = {beamer,colframe=blue!50!black,colback=blue!50,colupper=yellow!50},
]
\posterbox{name=A,column=1,at=middle}{First box}
\posterbox{name=B,column=2,at=row1}{Second box}
\end{tcbposter}
\end{exdispExample}
\end{posterlocTcbKey}


\begin{posterlocTcbKey}[][doc new=2017-07-03]{between}{=\meta{name1} and \meta{name2}}{无默认值，初始未设置}
  将文本框放置在文本框 \meta{name1} 之下、另一文本框 \meta{name2} 之上。
  \meta{name1} 与 \meta{name2} 也可以是预定义节点，
  见 \zvref{subsec:poster_environment}。
\begin{exdispExample}{between}
\begin{tcbposter}[
  poster = {showframe,height=3cm,spacing=2mm,rows=2},
  boxes  = {beamer,colframe=blue!50!black,colback=blue!50,colupper=yellow!50},
]
\posterbox{name=A,column=1,below=top}{First box}
\posterbox{name=B,column=1,between=A and bottom}{Second box}
\posterbox{name=C,column=2,above=bottom}{Third box}
\posterbox{name=D,column=2,between=top and C,span=2}{Fourth box}
\posterbox{name=E,column=3,between=D and bottom}{Fifth box}
\end{tcbposter}
\end{exdispExample}
\end{posterlocTcbKey}


\clearpage
\begin{posterlocTcbKey}[][doc new=2017-07-03]{sequence}{=\meta{sequence}}{无默认值，初始未设置}
  将当前文本框拆分为若干子文本框。这些子文本框按照给定的
  \meta{放置序列} 依次放置。\meta{序列} 的可用语法为：\par\medskip
  \meta{a 列} |between| \meta{名称 a1} |and| \meta{名称 a2} |then|\\
  \meta{b 列} |between| \meta{名称 b1} |and| \meta{名称 b2} |then|\\
  \meta{c 列} |between| \meta{名称 c1} |and| \meta{名称 c2} |then|\ldots\par\medskip
  显然，这会将第一个子文本框放置在 \meta{a 列} 中、
  \meta{名称 a1} 与 \meta{名称 a2} 之间；
  第二个子文本框放置在 \meta{b 列} 中、
  \meta{名称 b1} 与 \meta{名称 b2} 之间，依此类推。

\begin{exdispExample}{sequence}
\begin{tcbposter}[
  poster = {showframe,height=6cm,spacing=2mm,rows=2},
  boxes  = {beamer,colframe=blue!50!black,colback=blue!50,colupper=yellow!50},
]
\posterbox[adjusted title=A]{name=A,column=1,below=top,span=2}{First box}
\posterbox{name=B,column=2,above=bottom,span=2}{Second box}
\posterbox[adjusted title=C,colframe=red!50!black,colback=red!50]{
  name=C, sequence=1 between A and bottom then
                   2 between A and B then
                   3 between top and B
  }{\lipsum[2]}
\end{tcbposter}
\end{exdispExample}
\end{posterlocTcbKey}

\clearpage

\begin{docTcbKey}[][doc new=2017-07-03]{placeholder}{}{样式，无取值}
  若 \refKey{/tcb/posterloc/sequence} 的文本框内容太短，
  无法填满所有预留的子框，多余的空框将以 \refKey{/tcb/placeholder} 样式绘制。
  如果你希望空框什么都不绘制，可将此样式重定义为
  \refKey{/tcb/blankest}，例如。
\begin{exdispExample}{placeholder}
\begin{tcbposter}[
  poster = {showframe,height=2.5cm,spacing=2mm,rows=2},
  boxes  = {beamer,colframe=blue!50!black,colback=blue!50,colupper=yellow!50},
]
\posterbox{name=A,column=1,below=top,span=2}{First box}
\posterbox[colframe=red!50!black,colback=red!50]{
  name=B, sequence=1 between A and bottom then
                   2 between A and bottom then
                   3 between top and bottom
  }{Second box followed by placeholder boxes}
\end{tcbposter}
\end{exdispExample}
\end{docTcbKey}



\begin{posterlocTcbKey}[][doc new=2017-07-03]{xshift}{=\meta{length}}{无默认值，初始 |0pt|}
  将文本框水平移动 \meta{长度}。
\begin{exdispExample}{xshift}
\begin{tcbposter}[
  poster = {showframe,height=3cm,spacing=2mm,rows=2},
  boxes  = {beamer,colframe=blue!50!black,colback=blue!50,colupper=yellow!50},
]
\posterbox{name=A,column=1,row=1,xshift=6mm}{First box}
\posterbox{name=B,column=2,row=2,xshift=-6mm}{Second box}
\end{tcbposter}
\end{exdispExample}
\end{posterlocTcbKey}

\clearpage

\begin{posterlocTcbKey}[][doc new=2017-07-03]{yshift}{=\meta{length}}{无默认值，初始 |0pt|}
  将文本框垂直移动 \meta{长度}。
\begin{exdispExample}{yshift}
\begin{tcbposter}[
  poster = {showframe,height=3cm,spacing=2mm,rows=2},
  boxes  = {beamer,colframe=blue!50!black,colback=blue!50,colupper=yellow!50},
]
\posterbox{name=A,column=1,row=1,yshift=-4mm}{First box}
\posterbox{name=B,column=2,row=2,yshift=4mm}{Second box}
\end{tcbposter}
\end{exdispExample}
\end{posterlocTcbKey}


